% ===========================
% Columbia MSBA Assignment Template
% Author: Alexandre Sepúlveda de Dietrich
% ===========================
\documentclass[11pt]{article}

% --- Page & typography ---
\usepackage[margin=1in]{geometry}
\usepackage{microtype}
\usepackage{setspace}
\setstretch{1.08}

% --- Math, tables, figures ---
\usepackage{amsmath,amssymb}
\usepackage{siunitx}
\usepackage{booktabs, makecell}
\usepackage{graphicx}
\usepackage{caption}
\usepackage{subcaption}
\usepackage{float}
\usepackage{tikz}

% scale: 1 ft = 0.55 cm
\newcommand{\ftscale}{0.55}
\newcommand{\boardheight}{0.45}


% --- Links & smart refs ---
\usepackage{hyperref}
\hypersetup{
  colorlinks=true,
  linkcolor=blue!55!black,
  urlcolor=blue!55!black,
  citecolor=blue!55!black,
  pdfauthor={Alexandre Sepúlveda de Dietrich},
  pdftitle={MSBA Assignment}
}
\usepackage[nameinlink,capitalize]{cleveref}

% --- Header / footer ---
\usepackage{fancyhdr}

% --- Code listings (Python default); keep English identifiers/comments ---
\usepackage{listings}
\usepackage{xcolor}
\usepackage{inconsolata}

\usepackage[table]{xcolor}
\usepackage{multirow}

% ==========================================================================
% ==========================================================================
% ---------------------------
% CONFIG MACROS (edit per assignment)
% ---------------------------
\newcommand{\coursecode}{IEOR E4004}
\newcommand{\coursetitle}{Optimization Models and Methods}
\newcommand{\assignmenttitle}{Project I --- Childcare Desert Reduction in New York City}
\newcommand{\term}{Spring~2026}
\newcommand{\instructor}{Prof.\ Dr. Yaren Bilge Kaya}
\newcommand{\studentname}{Group 8}
\newcommand{\uni}{aes2399}
\newcommand{\studentemail}{aes2399@columbia.edu}
\newcommand{\dueday}{March 2026}
\newcommand{\logoimg}{}

% Toggle: show table of contents (0/1)
\newif\ifshowtoc
\showtoctrue
% Toggle: include Academic Integrity & AI Use section (0/1)
\newif\ifshowai
\showaitrue
% ==========================================================================
% ==========================================================================

% ---------------------------
% Styles & utilities
% ---------------------------
% siunitx numbers & alignment


% Regression helpers
\renewcommand{\arraystretch}{1.15}
\newcommand{\coef}[2]{#1\\ \scriptsize{(#2)}}
\newcommand{\ci}[2]{\shortstack[r]{[\,#1;\,#2\,]}}

% Math shortcuts
\newcommand{\E}{\mathbb{E}}
\newcommand{\Var}{\mathrm{Var}}
\newcommand{\Cov}{\mathrm{Cov}}
\newcommand{\1}{\mathbb{I}}

% LP formatting helpers
\newcommand{\lp}[3]{% #1 = max/min, #2 = objective, #3 = constraints (one per line with \\)
\[
\begin{array}{ll}
\text{#1} & #2 \\
\text{s.t.} & \begin{aligned}
#3
\end{aligned}
\end{array}
\]
}

% Optional: same but boxed (nice for examples)
\newcommand{\lpexpr}[3]{%
\begin{resultbox}
\lp{#1}{#2}{#3}
\end{resultbox}
}


% Graphics default path
\graphicspath{{figures/}}

% Listings: Python style
\definecolor{codebg}{HTML}{F8F8F8}
\definecolor{codeframe}{HTML}{E1E4E8}
\definecolor{codekw}{HTML}{005CC5}
\definecolor{codestr}{HTML}{22863A}
\definecolor{codecom}{HTML}{6A737D}
\definecolor{codegray}{HTML}{6E6E6E}
\lstdefinestyle{pycode}{
  language=Python,
  basicstyle=\linespread{0.98}\ttfamily\footnotesize,
  keywordstyle=\color{codekw}\bfseries,
  stringstyle=\color{codestr},
  commentstyle=\color{codecom}\itshape,
  numbers=left,
  numberstyle=\tiny\color{codegray},
  numbersep=8pt,
  frame=single,
  rulecolor=\color{codeframe},
  framerule=0.3pt,
  backgroundcolor=\color{codebg},
  showstringspaces=false,
  upquote=true,
  tabsize=2,
  keepspaces=true,
  columns=fullflexible,
  breaklines=true,
  breakatwhitespace=true,
  postbreak=\mbox{\textcolor{codegray}{$\hookrightarrow$}\space},
  xleftmargin=1.6em,
  framexleftmargin=1.6em,
  aboveskip=0.8\baselineskip,
  belowskip=0.8\baselineskip,
}
\lstset{style=pycode}

% Header/footer styles
\fancypagestyle{coverstyle}{%
  \fancyhf{}
  \lhead{Columbia University}
  \chead{\coursecode}
  \rhead{\studentname}
  \cfoot{\thepage}
  \renewcommand{\headrulewidth}{0.5pt}
  \renewcommand{\footrulewidth}{0.5pt}
}
\pagestyle{fancy}
\fancyhf{}
\lhead{\coursecode}
\rhead{\assignmenttitle}
\cfoot{\thepage}
\setlength{\headheight}{14pt}
\renewcommand{\headrulewidth}{0.5pt}
\renewcommand{\footrulewidth}{0.5pt}

% A lightweight "Problem" environment scaffold
\newenvironment{problem}[2][] % usage: \begin{problem}[Short name]{Problem N — Title}
{\section{#2 \ifx\relax#1\relax\else\ (\textit{#1})\fi}}
{\vspace{0.6em}}

% A small boxed “Result” callout
\usepackage{mdframed}
\newmdenv[
  linewidth=0.5pt,
  roundcorner=4pt,
  backgroundcolor=black!2,
  skipabove=\baselineskip,
  skipbelow=\baselineskip,
  innertopmargin=6pt,
  innerbottommargin=6pt,
  innerleftmargin=9pt,
  innerrightmargin=9pt
]{resultbox}

% ===========================
\begin{document}
% ===========================

% --- Title page ---
\begin{titlepage}
  \thispagestyle{coverstyle}
  \centering
  \vspace*{0.2cm}
  \vspace{0.5cm}
  {\Large \textsc{\coursecode{} — \coursetitle{}} \par}
  \vspace{0.2cm}
  {\LARGE \textbf{\assignmenttitle} \par}
  \vspace{0.3cm}
  {\large \term \par}
  \vspace{0.6cm}
  \noindent\makebox[\linewidth]{\rule{\linewidth}{1.2pt}}

    \vspace{0.9cm}
    \begin{minipage}{0.48\textwidth}
      \raggedright
      \textit{Team}\\
      Alexandre Sep\'ulveda de Dietrich (aes2399)\\
      Rong Ji\\
      Suyi Li\\
      Zichen He\\
      Zipeng Liu
    \end{minipage}\hfill
    \begin{minipage}{0.48\textwidth}
      \raggedleft
      \textit{Instructor}\\
      \instructor\\[4pt]
      \textit{Due}\\
      \dueday
    \end{minipage}

  \vfill
  {\large Department of Industrial Engineering and Operations Research\\Columbia University}
\end{titlepage}

\setcounter{page}{2}

\ifshowtoc
\tableofcontents
\newpage
\fi

% --- Academic Integrity & AI use disclosure ---
\ifshowai
\section*{Academic integrity \& AI use}
AI tools were used only as a writing and formatting assistant. The modeling approach, data preparation, optimization logic, interpretation of results, and final submitted content were produced, reviewed, and validated by the project team. AI assistance was limited to presentation tasks such as \LaTeX{} polishing, wording refinement, and formatting.
\addcontentsline{toc}{section}{Academic integrity \& AI use}
\fi

% --- Executive summary (optional but recommended for longer cases) ---
% \section*{Executive summary}
% In 6–10 sentences, state the problem, data, approach, principal quantitative results (use \num{} or math mode), limitations, and the final recommendation.
% Example numbers in math mode: $R^2=\num{0.613}$, expected margin $\approx \num{44.9}$ dollars.

% ===========================
% PROBLEM BLOCKS (duplicate as needed)
% ===========================
% ===========================
% PROJECT BODY
% ===========================

\section*{Project overview}
\addcontentsline{toc}{section}{Project overview}


\paragraph{Objective.}
This project studies how New York City can reduce childcare deserts at minimum cost by combining expansions of existing childcare facilities with the opening of new facilities. A zipcode is considered underserved when available childcare capacity is too low relative to the local child population. The policy challenge is therefore not only to increase total capacity, but also to ensure that enough capacity is available for children under age \(5\), while accounting for the financial and operational costs of intervention.

\paragraph{Project structure.}
We solve the problem in two stages. In \textbf{Part 1}, we formulate an ideal benchmark model in which the city can expand existing facilities and build new facilities freely within zipcodes. In \textbf{Part 2}, we extend this benchmark to a more realistic planning model by imposing stricter expansion limits, piecewise expansion costs, candidate-site restrictions, and minimum-distance requirements between facilities. This two-step structure is useful because it separates the core allocation logic from the implementation frictions that arise in practice.

\paragraph{Approach.}
The analysis combines data preparation, parameter estimation, and optimization. We first construct zipcode-level demand targets and baseline capacity measures from the provided demographic and childcare datasets. We then formulate a mixed-integer optimization model that chooses the least-cost combination of capacity expansions and new openings needed to eliminate childcare deserts. The ideal model in Part 1 provides a benchmark intervention plan, while the realistic model in Part 2 shows how that plan changes once siting and spacing constraints are introduced.

\paragraph{Managerial interpretation.}
From a policy perspective, the problem is a trade-off between two distinct levers. Expanding existing facilities is often cheaper because infrastructure is already in place, but expansion opportunities are limited. New construction is more expensive, but it may be necessary in zipcodes with large shortages or limited existing supply. The optimization model identifies the least-cost mix of these two interventions while meeting both total-capacity and under-\(5\) service requirements.

\section{Part 1: Ideal childcare expansion model}

\subsection{Problem statement}

In the first part of the project, we consider an idealized version of the childcare planning problem. The city is allowed to expand existing facilities and open new childcare centers in any zipcode, without explicit spatial siting restrictions. The objective is to eliminate all childcare deserts at minimum total cost.

A zipcode must satisfy two service requirements in order to be considered adequately covered. First, it must have enough \emph{total childcare capacity}. Second, it must have enough \emph{under-\(5\) childcare capacity}. The total-capacity requirement depends on local socioeconomic conditions. Following the project statement, a zipcode is treated as \emph{high demand} if either its employment rate is at least \(60\%\) or its average income is at most \(\$60{,}000\). In such zipcodes, total capacity must exceed one half of the child population. In all other zipcodes, total capacity must exceed one third of the child population. In every zipcode, under-\(5\) capacity must exceed two thirds of the local population aged \(0\) to \(5\).

This first model deliberately abstracts away from implementation frictions such as land availability and minimum-distance rules. Its purpose is to establish a benchmark: if the city could intervene in the most flexible way possible, what is the least-cost plan required to eliminate all childcare deserts?

\subsection{Data and parameter construction}

\subsubsection{Data cleaning}

Our analysis combines five raw datasets provided for the project. Before estimating model inputs, we first built a consistent data-cleaning pipeline to ensure that all datasets could be merged and used reliably.

We began by standardizing the structure of each file. Unnamed index columns were removed, zipcodes were reformatted into a uniform five-digit string, and variable names were harmonized across datasets. We also converted all relevant quantitative fields, such as capacities, population counts, employment rates, incomes, and coordinates, into numeric format so that later calculations would be based on consistent data types.

Next, we examined missing values and data quality issues. The main missingness appeared in the childcare facility dataset, where some facilities lacked latitude and longitude information. Rather than dropping these records, we retained them and created a flag to indicate missing geographic data, since coordinates are not essential for the ideal model but may matter in the realistic model. After merging the zipcode-level tables, missing employment and income values were imputed with citywide medians, while missing population entries were filled with zero so that the full zipcode universe remained in the planning dataset. We also created basic quality-control flags, including indicators for inactive facilities, zero total capacity, and any inconsistency between \texttt{children\_capacity} and total capacity.

To prepare the supply-side data, we filtered the childcare dataset to active facilities with positive capacity and used \texttt{children\_capacity} as a proxy for current under-\(5\) capacity in the cleaned facility data. On the demand side, we aggregated the population data to create zipcode-level measures of the child population aged \(0\) to \(12\) and the population aged \(0\) to \(5\), and cleaned the employment and income files into zipcode-level lookup tables. The potential locations dataset was also cleaned and assigned unique candidate IDs.

Finally, we constructed a master zipcode table by taking the union of zipcodes from all datasets. This ensured that the merged planning dataset covered the full set of relevant zipcodes rather than only the overlapping subset appearing in every source. As a result, the cleaned data provided a consistent foundation for parameter estimation and optimization.

\subsubsection{Demand-side parameters}

After cleaning the data, we translated the project requirements into model-ready demand parameters.

At the zipcode level, a zipcode \(z\) was classified as high demand if its employment rate was at least \(60\%\) or its average income was at most \(\$60{,}000\). Based on this rule, the required total childcare capacity was defined as one half of the child population aged \(0\) to \(12\) in high-demand zipcodes and one third otherwise. The required under-\(5\) capacity was defined as two thirds of the population aged \(0\) to \(5\). To implement these rules in the optimization model, we converted them into integer requirements:
\[
\mathrm{req\_total}_z =
\begin{cases}
\left\lfloor \dfrac{1}{2}\,\mathrm{pop\_0\_12}_z \right\rfloor + 1, & \text{if zipcode } z \text{ is high demand}\\[1.2ex]
\left\lfloor \dfrac{1}{3}\,\mathrm{pop\_0\_12}_z \right\rfloor + 1, & \text{otherwise}
\end{cases}
\]
and
\[
\mathrm{req\_under5}_z = \left\lceil \dfrac{2}{3}\,\mathrm{pop\_0\_5}_z \right\rceil
\]

This formulation reflects the project definition that a zipcode remains a desert if capacity is less than or equal to the threshold, so the requirement must be strictly exceeded.

For descriptive purposes, we also measured unmet need using
\[
\begin{aligned}
\mathrm{gap\_total}_z &= \max\!\left(0,\ \mathrm{req\_total}_z - \mathrm{existing\_total}_z\right),\\
\mathrm{gap\_under5}_z &= \max\!\left(0,\ \mathrm{req\_under5}_z - \mathrm{existing\_under5}_z\right)
\end{aligned}
\]

These two gap measures summarize the magnitude of the shortage in each zipcode and are used as diagnostic quantities in the analysis, even though the optimization model itself is written directly in terms of required and existing capacity.

\subsubsection{Supply-side parameters}

Existing supply was estimated by aggregating facility-level capacity within each zipcode. Existing total capacity was computed from current licensed capacity, and existing under-\(5\) capacity was approximated using \texttt{children\_capacity}. These aggregated quantities define the baseline level of service already available before any intervention.

At the facility level, for each active facility \(f\), we recorded current total capacity \(n_f\) and under-\(5\) proxy capacity. These values determine both the current supply and the feasible expansion opportunities of existing facilities.

\subsubsection{Expansion cost parameters}

The project statement allows expansion up to \(120\%\) of current capacity, subject to a maximum of \(500\) added slots. We interpret this rule as a bound on \emph{added slots}, so that
\[
x_f \le \min(1.2\,n_f,\ 500),
\]
where \(x_f\) denotes the number of added slots and
\[
U_f = \min(1.2\,n_f,\ 500)
\]
is the facility-specific upper bound. This interpretation is consistent with the assignment’s cost function, which explicitly distinguishes the cases \(x_f < n_f\) and \(x_f \ge n_f\), and therefore allows expansions that can exceed the facility’s original size. 

The project statement specifies that the first-block per-slot expansion cost \(c_{\text{slots}}(n_f)\) should decrease with current facility size, but it does not provide a closed-form expression. We therefore approximated the first expansion block with size-based per-slot cost tiers for expansions below \(100\%\) of current capacity:
\[
c_f^{(1)} =
\begin{cases}
350, & \text{if } n_f < 25\\
300, & \text{if } 25 \le n_f < 75\\
250, & \text{if } n_f \ge 75
\end{cases}
\]

This specification operationalizes the assignment's statement that the per-slot expansion cost \(c_{\text{slots}}(n_f)\) is decreasing in current facility size while keeping the implementation transparent.

For expansions beyond \(100\%\) of current capacity, we applied the prorated project formula
\[
c_f^{(2)} = \frac{20000 + 200\,n_f}{n_f}
\]

In the implementation, expansion was split into two decision blocks: one up to the current size and one beyond it, with bounds
\[
x^{(1)}_{f,\max} = \min(n_f, U_f),
\qquad
x^{(2)}_{f,\max} = \max(0, U_f - n_f)
\]

\subsubsection{New-facility parameters}

We included the three new-facility options specified in the project: small, medium, and large. Their characteristics are respectively \((100,50,\$65{,}000)\), \((200,100,\$95{,}000)\), and \((400,200,\$115{,}000)\), where each tuple denotes total capacity, under-\(5\) capacity, and fixed construction cost. We also incorporated the additional equipment cost of \(\$100\) for every newly created under-\(5\) slot. Together, these estimated parameters provide the direct inputs for the ideal optimization model. 

\subsection{Optimization model}

\subsubsection{Modeling assumptions}

To keep the benchmark model transparent, we adopt the following assumptions:

\begin{enumerate}
    \item Existing facilities may be expanded up to the maximum allowed by the assignment, subject to the facility-level cap on added slots.
    \item New facilities may be opened in any zipcode.
    \item New facilities come from a discrete menu of sizes, each with a fixed construction cost and a known contribution to total and under-\(5\) capacity.
    \item Demand is enforced at the zipcode level. The model does not represent cross-zipcode commuting or household travel behavior.
    \item Expansion and construction costs follow the project rules, with the unspecified first-block function \(c_{\text{slots}}(n_f)\) approximated by a decreasing three-tier schedule.
    \item The variable \texttt{children\_capacity} is used as a proxy for current under-\(5\) capacity at existing facilities.
    \item Higher expansion levels are represented through a separate decision block that uses the prorated cost coefficient from the project statement.
    \item Under-\(5\) capacity created through expansion and new construction is assigned explicitly through decision variables rather than assumed to follow mechanically from total capacity.
    \item Newly built facilities are not allowed to be expanded within the same planning horizon. 
\end{enumerate}

These assumptions make Part 1 a benchmark planning model rather than a fully operational deployment model. This is intentional: the goal is to first understand the core resource-allocation problem before introducing spatial and implementation constraints in Part 2.

\subsubsection{Decision variables}

Using the cleaned data and estimated parameters, we formulated an integer linear optimization model to determine the minimum funding needed to eliminate childcare deserts across New York City under the project's idealized assumptions.

\noindent
For each existing facility \(f\), we define:
\begin{itemize}
    \item \(x_{1f}\): added slots in the first expansion block, up to the facility’s current size;
    \item \(x_{2f}\): added slots in the second expansion block, for expansion beyond \(100\%\) of current size up to the project limit;
    \item \(u_f\): number of under-\(5\) slots assigned within the total expansion at facility \(f\).
\end{itemize}

\noindent
For each zipcode \(z\) and facility size \(s \in S=\{\text{small},\text{medium},\text{large}\}\), we define:
\begin{itemize}
    \item \(y_{zs}\): number of new facilities of size \(s\) built in zipcode \(z\);
    \item \(g_{zs}\): number of under-\(5\) slots assigned within those new facilities.
\end{itemize}

\noindent
All decision variables are restricted to be nonnegative integers.

\subsubsection{Objective function}

The objective is to minimize total funding, which includes three components:
\begin{enumerate}
    \item expansion costs at existing facilities;
    \item fixed construction costs for new facilities;
    \item additional equipment costs of \(\$100\) for each newly created under-\(5\) slot.
\end{enumerate}

\noindent
The objective can be written as
\[
\min \quad
\sum_{f \in F} \left(c_f^{(1)} x_{1f} + c_f^{(2)} x_{2f}\right)
+
\sum_{z \in Z} \sum_{s \in S} K_s y_{zs}
+
100\left(\sum_{f \in F} u_f + \sum_{z \in Z} \sum_{s \in S} g_{zs}\right)
\]

\subsubsection{Constraints}

The model includes four main groups of constraints.

First, each existing facility can expand only within its allowable limits:
\[
x_{1f} \le x^{(1)}_{f,\max},
\qquad
x_{2f} \le x^{(2)}_{f,\max}
\qquad \forall f \in F
\]

Second, the number of under-\(5\) slots added through expansion cannot exceed total added slots at the same facility:
\[
u_f \le x_{1f} + x_{2f}
\qquad \forall f \in F
\]

Third, each zipcode must satisfy the required total capacity after accounting for existing supply, expansions, and new construction:
\[
\mathrm{existing\_total}_z
+
\sum_{f \in F(z)} \left(x_{1f}+x_{2f}\right)
+
\sum_{s \in S} B_s y_{zs}
\ge
\mathrm{req\_total}_z
\qquad \forall z \in Z,
\]
where \(B_s\) is the total capacity of a new facility of size \(s\), and \(F(z)\) is the set of existing facilities located in zipcode \(z\).

Fourth, each zipcode must also satisfy the required under-\(5\) capacity:
\[
\mathrm{existing\_under5}_z
+
\sum_{f \in F(z)} u_f
+
\sum_{s \in S} g_{zs}
\ge
\mathrm{req\_under5}_z
\qquad \forall z \in Z
\]

Finally, under-\(5\) assignments in newly built facilities cannot exceed the under-\(5\) capacity allowed by each facility type:
\[
g_{zs} \le U_s\, y_{zs}
\qquad \forall z \in Z,\ \forall s \in S,
\]
where \(U_s\) is the maximum number of under-\(5\) slots available in a new facility of size \(s\).

All decision variables are integer and nonnegative:
\[
x_{1f}, x_{2f}, u_f \in \mathbb{Z}_+
\qquad \forall f \in F,
\]
\[
y_{zs}, g_{zs} \in \mathbb{Z}_+
\qquad \forall z \in Z,\ \forall s \in S
\]

Taken together, these constraints ensure that every zipcode receives enough total capacity and enough under-\(5\) capacity while respecting the project’s rules on expansion and construction.

\subsection{Results and interpretation}

\subsubsection{Optimal solution summary}

We solved the ideal optimization model using Gurobi. The solver returned status \texttt{OPTIMAL} in approximately \(0.43\) seconds and produced a feasible intervention plan that satisfies both the total-capacity and under-\(5\) requirements in every zipcode.

Before optimization, \(290\) zipcodes failed the total-capacity desert rule, \(179\) failed the under-\(5\) requirement, and \(308\) required intervention under at least one criterion. After optimization, all three counts fall to zero. Citywide capacity gaps are fully eliminated in the coded model: both the total-capacity gap and the under-\(5\) gap drop from positive pre-intervention levels to zero after implementation.

The optimal plan requires total funding of approximately \(\$132.81\) million. To achieve this, the model adds capacity through both expansion of existing facilities and construction of new facilities. Expansion contributes \(212{,}657\) additional total slots, including \(208{,}789\) under-\(5\) slots, indicating that existing facilities still play an important role in addressing shortages.

The model also builds \(487\) new facilities, adding \(154{,}300\) total slots and \(68{,}932\) under-\(5\) slots. The corresponding cost breakdown is \(\$55.74\) million for expansion, \(\$49.30\) million for new construction, and \(\$27.77\) million for under-\(5\) equipment. This highlights the importance of new construction in zipcodes with large unmet demand where expansion capacity is limited.

\begin{table}[H]
\centering
\small
\begin{tabular}{lS[table-format=9.0]}
\toprule
\textbf{Quantity} & {\textbf{Value}} \\
\midrule
Total intervention cost (\$) & 132800000 \\
Expanded total slots & 212657 \\
Expanded under-\(5\) slots & 208789 \\
Number of new facilities & 487 \\
New total slots & 154300 \\
New under-\(5\) slots & 68932 \\
\bottomrule
\end{tabular}
\caption{Summary of the optimal solution for the ideal model}
\end{table}

\begin{table}[H]
\centering
\small
\begin{tabular}{lrr}
\toprule
\textbf{Cost component} & \textbf{Amount (\$)} & \textbf{Share of total} \\
\midrule
Expansion cost & 55,740,297 & 41.97\% \\
New-build cost & 49,295,000 & 37.12\% \\
Under-\(5\) equipment cost & 27,772,100 & 20.91\% \\
\bottomrule
\end{tabular}
\caption{Cost breakdown for the ideal model}
\end{table}

\subsubsection{Interpretation}

The benchmark solution suggests that existing facilities are an important part of the policy response, but they cannot absorb the full shortage by themselves. New facilities remain necessary even in the ideal model, especially in the highest-intervention zipcodes such as \(11219\), \(11230\), \(11223\), \(11368\), and \(11220\), each of which requires several thousand additional slots. The large number of under-\(5\) slots added through expansion also suggests that the under-\(5\) requirement is a primary driver of the solution rather than a secondary constraint.

From a policy perspective, the ideal model shows that eliminating childcare deserts requires a large-scale intervention, not a marginal adjustment. Even before adding realistic siting restrictions, the city must combine widespread expansions with a substantial amount of new construction. This benchmark therefore provides a useful baseline estimate under the idealized assumptions used in Part 1.

\subsubsection{Limitations}

The strength of Part 1 is also its main limitation: it assumes a high degree of flexibility. In practice, new facilities cannot be opened arbitrarily within every zipcode, and existing facilities may face regulatory or physical barriers to expansion. The benchmark model also treats demand at the zipcode level and does not account for travel between neighborhoods. These limitations do not invalidate the model, but they do mean that its results should be interpreted as a planning benchmark rather than a deployable implementation plan. This motivates the realistic extension in Part 2.

\section{Part 2: Realistic childcare expansion model}

\subsection{Problem statement}

Part 2 keeps the same zipcode-level demand targets from Part 1, but it introduces three practical frictions that were omitted from the ideal benchmark. First, expansions of existing facilities are capped at only \(20\%\) of current capacity and priced through a three-block piecewise cost structure. Second, new facilities can no longer be opened arbitrarily within a zipcode; instead, each new opening must be assigned to an observed candidate site from the provided locations dataset. Third, candidate sites must satisfy a minimum-distance rule of \(0.06\) miles from one another, and candidate sites that are too close to existing facilities are excluded from consideration.

These additional restrictions make Part 2 a more operational planning model. Existing facilities remain fixed in place, but the city now has much less freedom to decide where new capacity can be added. The central question therefore changes from ``what is the cheapest way to eliminate childcare deserts under ideal flexibility?'' to ``how does the least-cost intervention plan change once expansion and siting become constrained?''

\subsection{Data and parameter construction}

\subsubsection{Realistic expansion parameters}

Part 2 reuses the cleaned shared datasets and the zipcode-level demand table from Part 1, but with one repair: zipcodes with zero observed child population are assigned zero total-slot requirement rather than a spurious requirement of one slot. For each active existing facility \(f\), we keep current total capacity \(n_f\) and the under-\(5\) proxy \(b_f\), then redefine the feasible expansion range to reflect the realistic policy recommendation that capacity cannot be increased by more than \(20\%\).

The realistic expansion cap is implemented through three blocks:
\[
\mathrm{cap1}_f = \left\lfloor 0.10\,n_f \right\rfloor,\qquad
\mathrm{cap2}_f = \left\lfloor 0.15\,n_f \right\rfloor - \left\lfloor 0.10\,n_f \right\rfloor,\qquad
\mathrm{cap3}_f = \left\lfloor 0.20\,n_f \right\rfloor - \left\lfloor 0.15\,n_f \right\rfloor
\]

These blocks correspond to the intervals \(0\%\) to \(10\%\), \(10\%\) to \(15\%\), and \(15\%\) to \(20\%\) of current capacity. Using cumulative floors preserves the intended total cap of \(\lfloor 0.20 n_f \rfloor\) at each facility. The repaired processed inputs imply a citywide realistic expansion capacity of \(61{,}197\) slots, which is still substantially smaller than the ideal benchmark's effective expansion room.

The corresponding block-specific cost coefficients follow the formulas used in the realistic parameter notebook:
\[
\mathrm{coef1}_f = \frac{20000 + 200\,n_f}{n_f},\qquad
\mathrm{coef2}_f = \frac{20000 + 400\,n_f}{n_f},\qquad
\mathrm{coef3}_f = \frac{20000 + 1000\,n_f}{n_f}
\]

\subsubsection{Candidate-site and distance parameters}

The realistic model uses the cleaned potential-locations dataset to define candidate sites for new construction. Each candidate site is combined with the same three build-size options from Part 1: small, medium, and large. The new-facility cost and capacity menu is therefore unchanged, but now every new opening must be placed at a specific observed site.

To operationalize the spacing requirement, we computed great-circle distances in miles between candidate sites and between candidate sites and existing facilities. A candidate site is blocked if it lies within \(0.06\) miles of an existing facility in the same zipcode. For candidate pairs in the same zipcode, any pair closer than \(0.06\) miles is added to a conflict table so that both sites cannot be selected simultaneously. Existing-existing pairs below the same threshold are tracked only diagnostically because existing facilities are fixed and cannot be removed.

\begin{table}[H]
\centering
\small
\begin{tabular}{lr}
\toprule
\textbf{Screening metric} & \textbf{Value} \\
\midrule
Candidate sites examined & 31,100 \\
Blocked candidate sites & 2,578 \\
Feasible candidate sites & 28,522 \\
Candidate-candidate conflict pairs & 11,455 \\
Candidate-existing conflict pairs & 4,730 \\
Existing-existing close pairs (diagnostic) & 7,644 \\
\bottomrule
\end{tabular}
\caption{Candidate-site screening summary for the realistic model}
\end{table}

\subsubsection{Build options and inherited assumptions}

The build menu from Part 1 is retained unchanged: small facilities contribute \((100, 50, \$65{,}000)\), medium facilities contribute \((200, 100, \$95{,}000)\), and large facilities contribute \((400, 200, \$115{,}000)\), where each tuple gives total capacity, under-\(5\) capacity, and fixed construction cost. As before, each newly created under-\(5\) slot carries an additional \(\$100\) equipment cost.

The demand side is mostly inherited from Part 1. Part 2 keeps the same high-demand classification rule, the same under-\(5\) requirement, and the same zipcode-level aggregation, but resets \(\mathrm{req\_total}_z\) to zero when the observed child population in zipcode \(z\) is zero. This removes \(131\) artificial total-demand requirements that otherwise force meaningless new builds in zero-demand areas.

\subsection{Optimization model}

\subsubsection{Modeling assumptions}

The realistic formulation uses the following assumptions:

\begin{enumerate}
    \item Demand requirements by zipcode follow Part 1 except that zero-child zipcodes are assigned zero total-slot requirement.
    \item Existing facilities remain fixed in place and may only expand within the realistic \(20\%\) cap.
    \item Expansion costs are modeled through three ordered blocks with increasing marginal cost.
    \item New facilities can only be opened at observed candidate sites.
    \item At most one facility size may be chosen at any candidate site.
    \item Candidate sites too close to existing facilities are blocked.
    \item Candidate sites that are too close to each other cannot both be selected.
    \item Existing-existing distance conflicts are recorded diagnostically but not enforced because the model does not relocate or close existing facilities.
    \item The same \texttt{children\_capacity} proxy is used for current under-\(5\) capacity at existing facilities.
\end{enumerate}

\subsubsection{Decision variables}

The realistic model extends the Part 1 decision structure by introducing candidate-site decisions. For each existing facility \(f\), we define:
\begin{itemize}
    \item \(x_{1f}, x_{2f}, x_{3f}\): expansion slots in the three realistic cost blocks;
    \item \(u_f\): under-\(5\) slots assigned within the expansion at facility \(f\).
\end{itemize}

\noindent
For each candidate site \(l\) and facility size \(s \in S\), we define:
\begin{itemize}
    \item \(y_{ls}\): a binary variable equal to \(1\) if a facility of size \(s\) is built at site \(l\);
    \item \(g_{ls}\): the number of under-\(5\) slots assigned within that selected new facility.
\end{itemize}

\subsubsection{Objective function}

The objective remains the minimization of total funding:
\[
\min \quad
\sum_{f \in F} \left(\mathrm{coef1}_f x_{1f} + \mathrm{coef2}_f x_{2f} + \mathrm{coef3}_f x_{3f}\right)
+
\sum_{l \in L} \sum_{s \in S} K_s y_{ls}
+
100\left(\sum_{f \in F} u_f + \sum_{l \in L} \sum_{s \in S} g_{ls}\right)
\]

As implemented in the repaired Part 2 pipeline, this objective uses additive block costs to preserve increasing marginal expansion costs in a tractable MILP. That is a modeling approximation to the assignment's more literal piecewise-total cost rule, which remains harder to solve at the same scale.

\subsubsection{Constraints}

The realistic model keeps the zipcode-level service constraints from Part 1 but adds several layers of feasibility restrictions.

First, each existing facility can expand only within the three block limits:
\[
x_{1f} \le \mathrm{cap1}_f,\qquad
x_{2f} \le \mathrm{cap2}_f,\qquad
x_{3f} \le \mathrm{cap3}_f
\qquad \forall f \in F
\]

Second, under-\(5\) slots created through expansion cannot exceed total added slots:
\[
u_f \le x_{1f} + x_{2f} + x_{3f}
\qquad \forall f \in F
\]

Third, each candidate site can host at most one facility size:
\[
\sum_{s \in S} y_{ls} \le 1
\qquad \forall l \in L
\]

Fourth, under-\(5\) assignments in selected new facilities cannot exceed the size-specific limit:
\[
g_{ls} \le U_s y_{ls}
\qquad \forall l \in L,\ \forall s \in S
\]

Fifth, blocked candidate sites are forced to remain unused, and any two candidate sites that form a conflict pair cannot both be selected:
\[
\sum_{s \in S} y_{ls} = 0
\qquad \forall l \in L_{\text{blocked}}
\]
\[
\sum_{s \in S} y_{l_1 s} + \sum_{s \in S} y_{l_2 s} \le 1
\qquad \forall (l_1,l_2) \in \mathcal{C}
\]

Finally, the zipcode-level total-capacity and under-\(5\)-capacity constraints are preserved from Part 1, except that new construction now enters through candidate sites rather than arbitrary zipcode-level build counts:
\[
\mathrm{existing\_total}_z
+
\sum_{f \in F(z)} (x_{1f}+x_{2f}+x_{3f})
+
\sum_{(l,s)\in L(z)} B_s y_{ls}
\ge \mathrm{req\_total}_z
\qquad \forall z \in Z
\]
\[
\mathrm{existing\_under5}_z
+
\sum_{f \in F(z)} u_f
+
\sum_{(l,s)\in L(z)} g_{ls}
\ge \mathrm{req\_under5}_z
\qquad \forall z \in Z
\]

\subsection{Results and interpretation}

\subsubsection{Optimal solution summary}

The repaired realistic mixed-integer model remains substantially larger than the ideal benchmark, with approximately \(201{,}980\) variables and \(155{,}819\) constraints. Gurobi solved the model in about \(42.9\) seconds and terminated within the prescribed \(0.1\%\) optimality tolerance, with a final relative MIP gap of roughly \(0.081\%\). The resulting objective value is approximately \(\$183.96\) million.

The cost composition still shifts sharply relative to Part 1, but the repaired Part 2 outputs are less wasteful than the earlier draft. Expansion accounts for \(\$9.38\) million of total spending, new facility construction accounts for \(\$146.81\) million, and the under-\(5\) equipment cost remains \(\$27.77\) million. The model no longer places any new facilities in zero-child zipcodes, and it uses the restored realistic expansion room before resorting to additional new construction.

\begin{table}[H]
\centering
\small
\begin{tabular}{lr}
\toprule
\textbf{Quantity} & \textbf{Value} \\
\midrule
Total intervention cost (\$) & 183,963,207 \\
Expansion cost (\$) & 9,381,107 \\
New-build cost (\$) & 146,810,000 \\
Under-\(5\) equipment cost (\$) & 27,772,100 \\
Expanded facilities & 1,126 \\
Expanded total slots & 25,108 \\
Expanded under-\(5\) slots & 25,067 \\
New facilities built & 1,284 \\
New total slots & 507,700 \\
New under-\(5\) slots & 252,654 \\
\bottomrule
\end{tabular}
\caption{Summary of the optimal solution for the realistic model}
\end{table}

The solution meets the coded total-capacity and under-\(5\) requirements in all \(311\) zipcodes. New construction still dominates the intervention mix, but less extremely than in the earlier draft: \(180\) zipcodes receive new facilities and \(151\) receive expansion. The size mix remains highly concentrated in the largest build option, with \(1{,}261\) large facilities, \(10\) medium facilities, and only \(13\) small facilities selected.

\begin{table}[H]
\centering
\small
\begin{tabular}{lrrr}
\toprule
\textbf{Metric} & \textbf{Part 1} & \textbf{Part 2} & \textbf{Change} \\
\midrule
Total cost (\$) & 132,807,397 & 183,963,207 & +51,155,810 \\
Expansion slots & 212,657 & 25,108 & -187,549 \\
New-build slots & 154,300 & 507,700 & +353,400 \\
New facilities & 487 & 1,284 & +797 \\
Under-\(5\) equipment cost (\$) & 27,772,100 & 27,772,100 & 0 \\
\bottomrule
\end{tabular}
\caption{Part 1 versus Part 2 intervention comparison}
\end{table}

\subsubsection{Interpretation}

Relative to the ideal benchmark, the realistic model still shifts the solution decisively away from expansion and toward new construction. The \(20\%\) cap on existing-facility growth leaves only limited room for expansion, and the minimum-distance rule reduces the menu of feasible build sites. Even after repairing the zero-demand zipcodes and restoring the full \(\lfloor 0.20 n_f \rfloor\) expansion cap, the city still needs \(1{,}284\) new facilities to deliver the same under-\(5\) coverage standard.

The under-\(5\) requirement remains the tightest system-wide constraint. In the final realistic solution, all \(311\) zipcodes have zero under-\(5\) slack, while \(130\) zipcodes have zero total-capacity slack. This still implies substantial total-capacity oversupply, but the repaired model wastes far fewer slots than the earlier Part 2 draft because it no longer builds into zero-demand areas.

The heaviest intervention remains concentrated in a handful of high-need zipcodes. The largest added-capacity totals occur in zipcodes \(11219\), \(11230\), \(11368\), \(11206\), \(11223\), and \(11204\), all of which require several thousand additional slots even after realistic restrictions are imposed. These neighborhoods therefore remain the clearest examples of where new construction is necessary rather than optional.

\subsubsection{Limitations}

Although Part 2 is more operational than the benchmark model, it is still not a full deployment model. First, existing-existing distance conflicts are recorded diagnostically but are not enforced because existing facilities are fixed and cannot be closed or relocated within the current planning horizon. Second, distance screening is enforced only within zipcodes, so the model does not represent cross-zipcode competitive effects or household travel across boundaries. Third, \(169\) existing facilities still lack coordinates, which means that candidate-existing screening is only as accurate as the available geographic data. Finally, the repaired Part 2 pipeline still uses an additive block-cost approximation rather than the literal piecewise-total realistic cost rule, because the exact version is materially slower to solve at the current citywide scale.

\section*{Conclusions}
\addcontentsline{toc}{section}{Conclusions}

The two-stage analysis provides a clear benchmark-to-implementation comparison. Part 1 shows that childcare deserts can be eliminated in the coded benchmark model for approximately \(\$132.81\) million when expansions and siting are highly flexible. The repaired Part 2 model shows that once realistic expansion caps and spacing restrictions are imposed, the required spending rises to approximately \(\$183.96\) million, an increase of about \(38.5\%\).

The comparison also clarifies how the intervention strategy changes. Under ideal flexibility, the solution relies heavily on expansion of existing facilities. Under realistic constraints, that option becomes far less important, and the city must instead finance a much larger program of new construction. Across both parts, however, the under-\(5\) requirement remains the central planning driver. This suggests that any practical policy discussion should focus not only on total slot creation, but also on where and how under-\(5\) capacity can be delivered at scale.

\section*{Bibliography}
\addcontentsline{toc}{section}{Bibliography}

\begin{thebibliography}{9}
\bibitem{assignment}
IEOR E4004. \emph{Project I: Eliminating Child Care Deserts in New York City through Optimization}. Columbia University, Spring 2026.

\bibitem{factsheet}
IEOR E4004. \emph{Project I: Eliminating Child Care Deserts in New York City through Optimization --- Fact Sheet}. Columbia University, Spring 2026.

\bibitem{americanprogress}
Center for American Progress. \emph{Child Care Deserts in the United States}. Available at \url{https://www.americanprogress.org/series/child-care-deserts/}.
\end{thebibliography}

% ===========================
% APPENDIX PLACEHOLDER
% ===========================
\appendix

\section*{Appendix}
\addcontentsline{toc}{section}{Appendix}

Detailed solver logs, zipcode-level outputs, facility-level outputs, and exported tables used to prepare the report are available in the project repository.

% ===========================
\end{document}
